%=======================================================================
%  CHAPTER 6 — ANOMALY / FRAUD DETECTION  (3 hrs)
%=======================================================================
\section{Anomaly / Fraud Detection}

{\color{sub}\footnotesize 3 hrs \gap$\cdot$\gap \mk{in 23/23 papers, the only chapter with perfect
coverage} \gap$\cdot$\gap syllabus gives it no sub-topics \gap$\cdot$\gap 1 numerical
$\rightarrow$ Ch 6 companion.}

\lead{Read this first:}
\begin{itemize}
  \item Only 3 hrs of syllabus, but it has appeared in \textbf{every single paper}.
        \mk{Best marks-per-hour in the subject.}
  \item One question dominates: \textbf{``What is anomaly detection $+$ explain distance-based
        methods''} \tS{10}. Learn that and the importance/types answer, and the chapter is banked.
\end{itemize}

\hr

\T{6.1 Basics}

\Q{What is anomaly detection? Why is it important?}{\m{2+2+4}\m{3+3}\m{5}\m{2+3}}
{\color{sub}\footnotesize \tF{6} \yr{82 Bh, 81 Bh, 81 Ba, 79 Bh, 73 Ch, 72 Ka}}

\sbs{%
\lead{Definition}
\begin{itemize}
  \item An \textbf{outlier / anomaly} is a data object that \textbf{deviates significantly} from
        the rest, as if generated by a \textbf{different mechanism}.
  \item \textbf{Anomaly detection} $=$ finding those objects.
  \item Hawkins: an observation deviating so much that it arouses suspicion it was generated by a
        different process.
\end{itemize}
\lead{\mk{Outlier $\ne$ noise}}
\begin{itemize}
  \item \textbf{Noise} $=$ random error / variance in a measured value. Has \textbf{no meaning},
        should be \textbf{removed} before mining.
  \item \textbf{Outlier} $=$ a genuine object from a different mechanism. Carries
        \textbf{information}, and is often \emph{the thing you are looking for}.
  \item \mk{Same data point, opposite treatment} depending on the task.
  \item Related: \textbf{novelty detection} finds genuinely new patterns; once confirmed they are
        \emph{folded into} the normal model, so later instances are no longer flagged.
\end{itemize}}{%
\lead{Why it matters / applications}
\begin{itemize}
  \item \textbf{Fraud detection}: credit-card use unlike the customer's normal pattern.
  \item \textbf{Intrusion detection}: traffic unlike normal network behaviour.
  \item \textbf{Medical}: unusual symptoms or test results $\rightarrow$ disease.
  \item \textbf{Public health}: a spike in one region $\rightarrow$ outbreak.
  \item \textbf{Industry}: fault and defect detection, equipment failure.
  \item \textbf{Ecosystem disturbance}: hurricanes, floods, droughts.
  \item \textbf{Data cleaning}: find and remove genuinely erroneous values.
\end{itemize}
\lead{\mk{The one-line justification}}
\begin{itemize}
  \item In these applications the \textbf{rare event is the valuable event}. Ordinary mining is
        tuned to find \emph{frequent} patterns, so it actively discards exactly what matters here.
\end{itemize}}

\vspace{3pt}
\creamq{Types of outliers \mk{(80 Bh asks this directly, \m{5})}}{}

\sbs{%
\begin{itemize}
  \item \textbf{Global outlier (point anomaly)}
  \begin{itemize}
    \item Deviates significantly from \textbf{the rest of the whole data set}.
    \item Simplest type; most methods target these.
    \item Eg a computer broadcasting a huge number of packets in a short time.
  \end{itemize}
  \item \textbf{Contextual (conditional) outlier}
  \begin{itemize}
    \item Deviates significantly \textbf{with respect to a chosen context}.
    \item Attributes split into two groups:
    \begin{itemize}
      \item \textbf{Contextual attributes}: define the context. Eg date, location.
      \item \textbf{Behavioural attributes}: what is judged. Eg temperature, humidity.
    \end{itemize}
    \item Eg \mk{28$^\circ$C is an outlier in a Kathmandu winter, normal in summer.}
    \item Generalises \textbf{local outliers}; a global outlier is the special case where the
          context is empty.
  \end{itemize}
\end{itemize}}{%
\begin{itemize}
  \item \textbf{Collective outlier}
  \begin{itemize}
    \item A \textbf{subset of objects that as a whole} deviates, even though
          \mk{no individual member is an outlier}.
    \item Eg 1 delayed shipment is normal; \textbf{100 delayed on one day} is a collective outlier.
    \item Eg a burst of packets from many machines at once $=$ denial-of-service, though each
          single packet looks fine.
    \item Detecting them needs the \textbf{relationships} between objects (sequence, spatial,
          graph structure), not just the objects.
  \end{itemize}
\end{itemize}
\lead{Challenges / issues \yr{76 Ch \m{2}, 72 Ch short note}}
\begin{itemize}
  \item \textbf{Defining ``normal''} is hard, and the boundary is rarely sharp.
  \item Normal behaviour \textbf{evolves} $\rightarrow$ today's model is stale tomorrow.
  \item Application-specific: one definition does not transfer between domains.
  \item \textbf{Scarcity of labelled data} for training and validation.
  \item \textbf{Noise looks like outliers} $\rightarrow$ hard to separate.
  \item \mk{Interpretability}: you must justify \emph{why} something was flagged.
  \item Handling \textbf{high dimensionality} (distances lose meaning).
\end{itemize}}

\hr

\T{6.2 Categories of Anomaly Detection Schemes}

\Q{Briefly explain different types of anomaly detection schemes.}{\m{4}\m{6}}
{\color{sub}\footnotesize \yr{81 Ba, 78 Bh} \gap 78 Bh asks to \textbf{compare and contrast three methods}.}

\sbsr{0.44}{%
\lead{(A) By how much labelling is available}
\begin{itemize}
  \item \textbf{Supervised}
  \begin{itemize}
    \item Training set labelled both normal \emph{and} outlier $\rightarrow$ treat as a
          \textbf{classification} problem.
    \item Problem: \mk{severe class imbalance} (outliers are rare) $+$ the sample may not
          represent the outlier distribution.
    \item Fixes: oversample outliers, make up artificial outliers, cost-sensitive learning.
  \end{itemize}
  \item \textbf{Semi-supervised}
  \begin{itemize}
    \item Only the \textbf{normal} class is labelled (common in practice).
    \item Build a model of normality; anything not matching it is an outlier.
  \end{itemize}
  \item \textbf{Unsupervised}
  \begin{itemize}
    \item \textbf{No labels at all}. Assumes normal objects are \emph{far more frequent} and
          cluster together.
    \item Most practical, and the default assumption in exam answers.
  \end{itemize}
\end{itemize}}{%
\lead{(B) By the assumption made \mk{(this is the ``three methods'' answer for 78 Bh)}}
\begin{tabularx}{\linewidth}{@{}L{2.5cm}XX@{}}
\toprule
\hrow \thead{Method} & \thead{Assumption} & \thead{Weakness} \\
\textbf{Statistical} (model-based) &
Normal data follows a \textbf{stochastic model}; low-probability objects are outliers. &
Must guess the right distribution; poor in high dims. \\
\textbf{Proximity-based} &
An outlier is \textbf{far} from its neighbours (or in a \textbf{sparse} region). &
$O(n^2)$ cost; sensitive to $r$ and $k$; struggles w/ varying density (distance-based only). \\
\textbf{Clustering-based} &
Normal objects belong to \textbf{large, dense} clusters; outliers do not. &
Depends on the clustering; costly; outliers can distort the clusters themselves. \\
\bottomrule
\end{tabularx}
\begin{itemize}
  \item \textbf{Classification-based} is sometimes given as a 4th: train a 1-class model of normality
        (eg 1-class SVM), flag whatever falls outside.
\end{itemize}}

\hr

\T{6.3 Distance / Proximity-Based Approaches}

\Q{Explain distance-based methods for anomaly detection.}{\m{2+3}\m{8}\m{2+6}\m{10}\m{5}\m{4+3}}
{\color{sub}\footnotesize \tS{10} \yr{82 Bh, 81 Bh, 80 Ba, 79 Ba, 75 Ch, 75 Ash, 74 Ch, 73 Shr, 71 Ch, 70 Asa}
\gap \mk{The single most repeated question in this chapter.}}

\sbsr{0.52}{%
\lead{Core idea}
\begin{itemize}
  \item Judge an object by its \textbf{proximity to its neighbours}.
  \item An object is an outlier if its neighbourhood is \textbf{sparse} or its neighbours are \textbf{far}.
  \item Needs only a distance measure {\footnotesize$\rightarrow$ Ch 2, §2.4}, no distribution assumption.
\end{itemize}
\lead{1. Distance-based: the DB$(r,\pi)$ outlier}
\begin{itemize}
  \item Fix a \textbf{distance threshold $r$} (defines the neighbourhood) and a
        \textbf{fraction threshold $\pi$}.
  \item Object $o$ is a \textbf{DB$(r,\pi)$-outlier} if:
        \[ \frac{\lVert \{o' \mid dist(o,o') \le r\}\rVert}{\lVert D\rVert} \le \pi \]
  \item In words: \mk{fewer than a $\pi$ fraction of all objects lie within distance $r$ of $o$.}
  \item \textbf{Equivalent $k$-NN form}: with $k = \lceil \pi\lVert D\rVert \rceil$, object $o$ is
        an outlier if $dist(o, o_k) > r$, where $o_k$ is its $k$th nearest neighbour.
  \item \textbf{Nested-loop algorithm}: for each object, scan the others and count how many fall
        within $r$; stop early once the count exceeds $\pi n$. Cost $O(n^2)$.
\end{itemize}
\lead{2. $k$-NN score (the other common phrasing)}
\begin{itemize}
  \item Score each object by the \textbf{distance to its $k$th nearest neighbour}
        (or the average of its $k$ nearest).
  \item Rank all objects by that score; the top ones are the outliers.
  \item \yr{71 Shr asks exactly this, \m{10}.}
\end{itemize}}{%
\lead{3. Density-based: the local outlier \yr{80 Bh \m{3}}}
\begin{itemize}
  \item \mk{Why needed}: distance-based methods take a \textbf{global} view, using one $r$ for the
        whole data set.
  \item If the data has \textbf{one dense cluster and one sparse cluster}, no single $r$ works:
  \begin{itemize}
    \item Set $r$ small $\rightarrow$ the entire sparse cluster is wrongly flagged.
    \item Set $r$ large $\rightarrow$ genuine outliers beside the dense cluster are missed.
  \end{itemize}
  \item \textbf{Fix}: compare an object's density to \textbf{the density of its own neighbours}, not
        to the whole data set.
  \item \textbf{Local Outlier Factor (LOF)}:
  \begin{itemize}
    \item $LOF \approx 1$ $\rightarrow$ density like its neighbours $\rightarrow$ normal.
    \item $LOF \gg 1$ $\rightarrow$ much sparser than its neighbours $\rightarrow$ outlier.
    \item Built from \emph{reachability distance} and \emph{local reachability density}.
  \end{itemize}
  \item \mk{Distance-based finds global outliers; density-based finds local ones.} That contrast is
        the marks.
\end{itemize}
\lead{Strengths / weaknesses}
\begin{itemize}
  \item $+$ No distribution assumption; simple; works on any data w/ a distance measure.
  \item $-$ $O(n^2)$ in the naive form.
  \item $-$ Sensitive to the choice of $r$, $\pi$, $k$.
  \item $-$ Degrades in high dimensions (curse of dimensionality).
\end{itemize}}

\vspace{3pt}
\sbs{%
\creamq{Clustering-based \yr{76 Ash short note}}{}
\begin{itemize}
  \item Cluster the data, then flag as outliers the objects that:
  \begin{itemize}
    \item belong to \textbf{no cluster} (eg DBSCAN noise points),
    \item are \textbf{far from their own cluster centroid}, or
    \item belong to a cluster that is \textbf{very small or very sparse}.
  \end{itemize}
  \item \textbf{DBSCAN} gives it free: \textbf{noise points} are exactly the outliers.
        {\footnotesize$\rightarrow$ Ch 5}
  \item $+$ Unsupervised; reuses clustering you may already have.
  \item $-$ Effectiveness depends entirely on the clustering; outliers can themselves
        distort the clusters.
\end{itemize}}{%
\creamq{Statistical approaches}{}
\begin{itemize}
  \item Assume a model, flag low-probability objects.
  \item \textbf{Parametric}: assume a distribution (usually \textbf{Gaussian}), estimate its
        parameters, flag points in the tails.
  \begin{itemize}
    \item Common rule: more than \textbf{3 standard deviations} from the mean.
    \item \textbf{Grubbs' test}: detects one outlier at a time in univariate normal data, removes
          it, repeats. $H_0$: no outlier.
  \end{itemize}
  \item \textbf{Non-parametric}: no assumed shape; use histograms or kernel density estimates.
  \item \textbf{Convex hull method} \yr{70 Ch asks this by name} \added
  \begin{itemize}
    \item Points on the \textbf{convex hull} (the boundary of the smallest convex polygon
          enclosing all points) are assumed to be \textbf{extreme}, hence outliers.
    \item Peel the hull layer by layer to get successive outlier candidates.
    \item \mk{Limitation}: only ever finds \emph{boundary} points. An outlier sitting \emph{inside}
          the data cloud, in a sparse pocket, is never detected.
  \end{itemize}
\end{itemize}}

\penalty0\vspace{3pt}
\creamq{The likelihood (mixture-model) approach}{\m{8}}
{\color{sub}\footnotesize \yr{70 Asa} \gap ``What is anomaly detection? Explain likelihood approach
for anomaly detection.'' \gap Tan/Steinbach \S10.2; not in Han \& Kamber. \added}

\sbs{%
\lead{The model}
\begin{itemize}
  \item Treat the data set $D$ as a \textbf{mixture of two distributions}:
  \begin{itemize}
    \item $M$ $=$ the \textbf{majority} (normal) distribution, something tight such as a Gaussian.
    \item $A$ $=$ the \textbf{anomaly} distribution, taken as \textbf{uniform} because nothing is
          assumed about how outliers are spread.
  \end{itemize}
  \item $\lambda$ $=$ prior probability that an object is an anomaly, a small number fixed in
        advance.
  \[ D = (1-\lambda)\,M + \lambda\,A \]
  \item At any moment the objects are partitioned into $M_t$ and $A_t$.
        \mk{Start with $M_0 = D$ and $A_0 = \emptyset$}: assume everything is normal.
\end{itemize}

\lead{Log-likelihood of the whole data set}
\[ LL_t(D) = |M_t|\log(1-\lambda) + \sum_{x_i \in M_t}\log P_{M_t}(x_i) \]
\[ \qquad {} + |A_t|\log \lambda + \sum_{x_i \in A_t}\log P_{A_t}(x_i) \]
\begin{itemize}
  \item First two terms are the \emph{normal} account, last two the \emph{anomaly} account.
\end{itemize}}{%
\lead{Algorithm}
\begin{enumerate}
  \item $M \leftarrow D$, \; $A \leftarrow \emptyset$, compute $LL_t(D)$.
  \item For each point $x \in M$:
  \begin{enumerate}
    \item Tentatively move $x$ to $A$.
    \item Re-estimate $M$'s parameters without $x$ and compute the new $LL_{t+1}(D)$.
    \item $\Delta = LL_{t+1}(D) - LL_t(D)$.
    \item If $\Delta > c$ (a threshold), \mk{keep $x$ in $A$ and declare it an anomaly};
          otherwise put it back in $M$.
  \end{enumerate}
  \item Repeat until no point moves.
\end{enumerate}

\lead{Why the test works \mk{(this is the explanation the marks are for)}}
\begin{itemize}
  \item Moving a \emph{normal} point out forces $M$ to explain fewer points it already explained
        well, and the uniform $A$ explains it badly $\rightarrow$ likelihood \textbf{falls}.
  \item Moving a \emph{genuine outlier} out lets $M$ tighten around the rest, and the outlier was
        costing $M$ a huge penalty anyway $\rightarrow$ likelihood \textbf{rises sharply}.
  \item So \mk{a large gain in likelihood is itself the evidence} that the point does not belong.
\end{itemize}

\lead{Strengths and weaknesses}
\begin{itemize}
  \item $+$ Firm statistical footing; gives a principled test rather than an arbitrary distance.
  \item $+$ Handles mixed / categorical attributes if a model exists for them.
  \item $-$ Needs a distribution assumption for $M$, and a wrong one gives wrong outliers.
  \item $-$ Expensive: $M$ is re-estimated once per candidate point.
  \item $-$ The threshold $c$ and the prior $\lambda$ are both guesses.
  \item $-$ Degrades badly in \textbf{high dimensions}, like every parametric method.
\end{itemize}}

\hr

\T{6.4 Base Rate Fallacy}

\Q{Base rate fallacy}{\m{3}\m{5}}
{\color{sub}\footnotesize \yr{70 Ch short note, 69 Ch \m{7}, 76 Ch numerical} \gap
\mk{Numerical solved in the Ch 6 companion.} Not in Han \& Kamber; from Tan/Steinbach and Axelsson. \added}

\sbs{%
\lead{The fallacy}
\begin{itemize}
  \item \textbf{Base rate} $=$ the prior probability of the rare event, eg $P(\text{intrusion})$.
  \item \textbf{Base rate fallacy} $=$ the tendency to \textbf{ignore the base rate} and judge only
        by the test's accuracy.
  \item Consequence: when the event is very rare, \mk{even a highly accurate detector produces
        mostly false alarms}, because the huge normal population generates more false positives
        than the tiny anomalous population generates true positives.
\end{itemize}
\lead{Set-up (intrusion detection)}
\begin{itemize}
  \item $I$ $=$ intrusive behaviour, $\lnot I$ $=$ normal. \quad $A$ $=$ alarm, $\lnot A$ $=$ no alarm.
  \item \textbf{Detection rate} $= P(A|I)$ \quad (true positive rate)
  \item \textbf{False alarm rate} $= P(A|\lnot I)$
  \item \textbf{Bayesian detection rate} $= P(I|A)$ \mk{$\rightarrow$ what actually matters to the
        analyst: given an alarm, is it real?}
\end{itemize}}{%
\lead{Bayes}
\[ P(I|A) = \frac{P(A|I)\,P(I)}{P(A|I)P(I) + P(A|\lnot I)P(\lnot I)} \]
\lead{Worked illustration}
\begin{itemize}
  \item Test 99\% accurate both ways, disease prevalence 1 in 100.
  \item Then $P(\text{disease}\mid\text{positive})$ is only about \mk{1/2}, not 99\%.
  \item With a rarer event the result gets far worse.
\end{itemize}
\lead{Axelsson's conclusion}
\begin{itemize}
  \item In intrusion detection $P(I)$ is \textbf{tiny}, so the false alarm rate
        \textbf{dominates} the equation.
  \item $\rightarrow$ you need an \mk{extremely low false alarm rate} to get a usable Bayesian
        detection rate. Raw accuracy is not enough.
  \item This is the formal reason \textbf{accuracy is the wrong metric} for anomaly detection
        {\footnotesize$\rightarrow$ same argument as class imbalance in Ch 3, §3.7}.
\end{itemize}}

\hr

%-----------------------------------------------------------------------
\T{6.5 PYQ Mapping}

{\color{sub}\footnotesize \mk{N} $=$ solved in the Ch 6 Numerical companion.
The syllabus gives this chapter no sub-topics, so PYQs are grouped by what they ask.}

\begin{enumerate}
  \item \tS{6} \textbf{Anomaly detection $+$ distance-based methods.} \mk{The core question.}
  \begin{itemize}
    \item ``What is anomaly detection? Explain distance based method.'' \hfill \m{2+3} \yr{80 Ba}
    \item same wording \hfill \m{8} \yr{73 Shr}
    \item ``Why is it important? $+$ distance based method'' \hfill \m{2+6} \yr{75 Ash}
    \item ``What do you mean by anomaly detection and why is it important? Describe distance based approaches.'' \hfill \m{4+3} \yr{74 Ch}
    \item ``Explain few distance based approaches'' \hfill \m{10} \yr{71 Ch}
    \item ``What is outlier? Explain the distance base approaches.'' \hfill \m{5} \yr{75 Ch}
  \end{itemize}
  \item \tF{3} \textbf{Definition $+$ importance $+$ types of schemes.}
  \begin{itemize}
    \item ``What is it? Why important? Briefly explain different types of schemes.'' \hfill \m{2+2+4} \yr{81 Ba}
    \item ``Why is it important and where is it applicable?'' \hfill \m{3+3} \yr{79 Bh}
    \item ``Discuss its importance in security.'' \hfill \m{5} \yr{72 Ka}
  \end{itemize}
  \item \tP{1} Types of outlier w/ examples $+$ how density-based detection works. \hfill \m{5+3} \yr{80 Bh}
  \item \tP{1} \textbf{Compare and contrast 3 different methods} of anomaly detection. \hfill \m{6} \yr{78 Bh}
  \item \tP{1} Issues related to anomaly detection \hfill \m{2} \yr{76 Ch};
        \mk{N} false alarm rate / detection rate numerical \hfill \m{3} \yr{76 Ch}
  \item \tP{1} How can the \textbf{Nearest-Neighbor} algorithm be used for anomaly detection? \hfill \m{10} \yr{71 Shr}
  \item \tP{4} Short notes:
  \begin{itemize}
    \item Clustering and its application in anomaly detection \hfill \m{3} \yr{76 Ash}
    \item Issues in anomaly / fraud detection \hfill \m{4} \yr{72 Ch}
    \item Data Mining for Anomaly Detection \hfill \m{4} \yr{74 Ash}
    \item \textbf{Anomaly Detection; Base Rate Fallacy; Convex Hull Method} \hfill \m{15} \yr{70 Ch}
  \end{itemize}
\end{enumerate}

\creamq{Papers added by the 2026-09 merge}{}
{\color{sub}\footnotesize The six papers of \code{DM4.1BCT.pdf} (2082 Bh, 2081 Bh, 2079 Ba, 2073 Ch, 2070 Asa, 2069 Ch), indexed here for the
first time --- the notes were built before that scan was merged. \mk{N} $=$ solved in this chapter's companion. \added}
\begin{itemize}
  \item What is anomaly detection? $+$ distance-based method. \hfill \m{2+5} \yr{82 Bh}
  \item Anomaly detection $+$ the different \textbf{approaches}. \hfill \m{2+6} \yr{81 Bh}
  \item \textbf{Strengths and weaknesses} of the statistical, proximity-based, density-based and cluster-based approaches. \hfill \m{6} \yr{79 Ba}
  \item Anomaly detection $+$ the \textbf{issues} associated with it. \hfill \m{2+3} \yr{73 Ch}
  \item Anomaly detection $+$ the \textbf{likelihood approach}. \hfill \m{8} \yr{70 Asa} \added
  \item What is \textbf{Base Rate Fallacy}? Explain w/ example. \hfill \m{7} \yr{69 Ch}
\end{itemize}

\lead{Coverage check}
\begin{itemize}
  \item All Ch 6 PYQs covered. 1 numerical \yr{76 Ch} $\rightarrow$ companion.
  \item \textbf{Gaps filled by research} \added:
  \begin{itemize}
    \item \textbf{Convex hull method} \yr{70 Ch}: named in the Tan slides but never explained,
          and absent from Han \& Kamber. Written out above w/ its limitation.
    \item \textbf{Base rate fallacy}: the notes give the psychology definition but not the
          intrusion-detection form. Axelsson's argument added.
    \item \textbf{Grubbs' test} and the supervised/semi-supervised/unsupervised split.
  \end{itemize}
\end{itemize}

\lead{Sources used for this chapter}
\begin{itemize}
  \item \textbf{Han \& Kamber 3rd ed}, ch 12 Outlier Detection: definition, the 3 outlier types,
        method taxonomy, DB$(r,\pi)$ formula, density-based / LOF. \mk{Primary authority.}
  \item \code{Anamoly Detection/Anamoly Detection.pdf} (Tan, Steinbach \& Kumar slides):
        convex hull, Grubbs' test, base rate fallacy, Axelsson.
  \item \code{data mining\_/Chapter 6 \_ 7.pdf} (Bhupesh K. Mishra), \code{Intrusion.pptx},
        \textbf{DBK Lecture-17}.
\end{itemize}
